\documentclass[aspectratio=169]{beamer}
\usetheme{Madrid}
\usecolortheme{default}

\usepackage{amsmath}
\usepackage{amssymb}
\usepackage{graphicx}
\usepackage{booktabs}

\title{From Theory to Practice: My Journey with Learning Machines}
\subtitle{STAT 747 - Homework 3 Reflection}
\author{Student Reflection}
\date{\today}

\begin{document}

% Title slide
\begin{frame}
\titlepage
\end{frame}

% ============================================================================
% PART 1: THE PHILOSOPHICAL FOUNDATION (2 minutes)
% ============================================================================

\section{The Philosophical Foundation}

% Slide 1: Tree-based vs Kernel Methods
\begin{frame}{Two Philosophical Approaches to Complexity}
\begin{columns}[T]
\begin{column}{0.48\textwidth}
\textbf{Tree-Based Methods}
\begin{itemize}
\item Recursive partitioning
\item "Divide and conquer"
\item Sequential yes/no questions
\item Inherently interpretable
\end{itemize}
\vspace{0.5cm}
\centering
\textit{"Is pixel 378 > 0.5?"}\\
\textit{If yes → left, else → right}
\end{column}

\begin{column}{0.48\textwidth}
\textbf{Kernel Methods (SVMs)}
\begin{itemize}
\item Geometric transformation
\item "Representation learning"
\item Lift to higher dimensions
\item Find optimal viewpoint
\end{itemize}
\vspace{0.5cm}
\centering
\textit{RBF kernel: map to}\\
\textit{infinite-dimensional space}
\end{column}
\end{columns}

\vspace{0.5cm}
\centering
\textbf{One seeks optimal questions, the other seeks optimal viewpoints}
\end{frame}

% Slide 2: Bias-Variance Evolution
\begin{frame}{My Evolution Understanding Bias-Variance}

\textbf{Before:} Mechanical thinking
\begin{itemize}
\item Complex models → low bias, high variance
\item Simple models → high bias, low variance
\end{itemize}

\vspace{0.5cm}
\textbf{After:} System thinking

\begin{center}
\Large
\textcolor{blue}{Bias-variance is a \textbf{model-data system property},\\
not just a model property}
\end{center}

\vspace{0.3cm}
\begin{columns}[T]
\begin{column}{0.48\textwidth}
\textbf{Exercise 2 (Economic Data)}\\
Linear SVR "simple" but \textbf{optimal}\\
RMSE: 0.0087 (best)\\
\small{Linear data → linear model wins}
\end{column}

\begin{column}{0.48\textwidth}
\textbf{Exercise 3 (MNIST Digits)}\\
RBF SVM "complex" but \textbf{necessary}\\
Accuracy: 97.2\% (best)\\
\small{Curved boundaries → kernels win}
\end{column}
\end{columns}

\end{frame}

% Slide 3: Most Elegant Insight
\begin{frame}{Most Elegant Mathematical Insight}

\textbf{Alice's Broken Dual Formulation}

\vspace{0.3cm}
She removed: $\alpha_i \leq C$ (thought it was redundant)

\vspace{0.5cm}
\begin{center}
\Large
\textcolor{red}{Constraints encode problem structure}
\end{center}

\vspace{0.5cm}
\textbf{Without $\alpha_i \leq C$:}
\begin{itemize}
\item Dual becomes unbounded on non-separable data
\item Lagrange multipliers → $\infty$ as optimizer forces impossible separation
\item No mathematical or philosophical coherence
\end{itemize}

\vspace{0.3cm}
\textbf{With $\alpha_i \leq C$:}
\begin{itemize}
\item "No single training example should dominate model behavior"
\item Mathematically necessary + philosophically sensible
\item Perfect primal-dual correspondence
\end{itemize}

\end{frame}

% ============================================================================
% PART 2: PRACTICAL WISDOM (3 minutes)
% ============================================================================

\section{Practical Wisdom}

% Slide 4: Most Surprising Finding - Regression
\begin{frame}{Most Surprising Finding: Simplicity Wins}

\textbf{Exercise 2: Economic Growth Prediction (54 countries)}

\begin{center}
\begin{tabular}{lcc}
\toprule
\textbf{Model} & \textbf{Test RMSE} & \textbf{Training Time} \\
\midrule
Linear SVR & \textcolor{blue}{\textbf{0.0087}} & 0.014s \\
Polynomial SVR & 0.0102 & 0.018s \\
RBF SVR & 0.0105 & 0.021s \\
\midrule
Decision Tree & 0.0154 & 0.003s \\
Random Forest & 0.0138 & 0.142s \\
Gradient Boosting & 0.0123 & \textcolor{red}{0.247s} \\
\bottomrule
\end{tabular}
\end{center}

\vspace{0.3cm}
\textbf{The Surprise:} I engineered interaction terms and polynomial features, expecting ensemble methods to dominate.

\vspace{0.2cm}
\textbf{The Reality:} Feature engineering + linear model crushed "sophisticated" methods.

\vspace{0.2cm}
\centering
\Large
\textcolor{blue}{Thoughtful simplicity beats automated complexity}

\end{frame}

% Slide 5: Partial Dependence - The Paradox
\begin{frame}{The Paradox: Random Forest Saw Non-Linearity But Lost}

\begin{center}
\includegraphics[width=0.85\textwidth]{output/partial_dependence.png}
\end{center}

\vspace{-0.2cm}
\small
Random Forest revealed genuine non-linear patterns, yet linear SVR still won.\\
\textbf{Why?} Feature engineering moved non-linearity into the feature space.

\end{frame}

% Slide 6: Key Insights from MNIST
\begin{frame}{MNIST Classification: Three Key Insights}

\begin{enumerate}
\item \textbf{Kernel selection = data structure matching}
\begin{itemize}
\item Linear SVM: 94.1\% (can't represent curved digit boundaries)
\item RBF SVM: \textcolor{blue}{\textbf{97.2\%}} (matches digit topology)
\item An '8' is two loops → needs non-linear boundaries
\end{itemize}

\vspace{0.3cm}
\item \textbf{High dimensions favor margin maximization}
\begin{itemize}
\item 784 pixels → 50 PCA components
\item Trees: axis-aligned splits inefficient (DT: 92.2\%, RF: 91.2\%)
\item SVMs: global geometric optimization (RBF: 97.2\%)
\end{itemize}

\vspace{0.3cm}
\item \textbf{Interpretability includes uncertainty}
\begin{itemize}
\item Decision function confidence: 0.82 (errors) vs 2.15 (correct)
\item \textcolor{blue}{The model knew when it didn't know}
\end{itemize}
\end{enumerate}

\end{frame}

% Slide 7: Misclassified Examples
\begin{frame}{The Model Knew When It Didn't Know}

\begin{center}
\includegraphics[width=0.75\textwidth]{output/svm_misclassified.png}
\end{center}

\vspace{-0.2cm}
\small
42 errors out of 1,500 (2.8\% error rate)\\
Most are genuinely ambiguous cases where even humans would struggle.\\
\textbf{Key:} Model signaled low confidence on these boundary cases.

\end{frame}

% Slide 8: Practical Advice
\begin{frame}{Six Practical Lessons Learned}

\begin{enumerate}
\item \textbf{Start simple, complicate mindfully}
\begin{itemize}
\item Always fit a linear baseline first
\item Complexity should be justified by performance
\end{itemize}

\item \textbf{Validate with diagnostics, not just metrics}
\begin{itemize}
\item Initial R² = -19.5 (impossible!)
\item Residual plots caught preprocessing errors
\end{itemize}

\item \textbf{Document test set statistics}
\begin{itemize}
\item Record: $n$, $\bar{y}$, $\sigma$, RMSE, MAE, R²
\item Enable algebraic verification: $R^2 = 1 - \frac{\text{RMSE}^2}{\sigma^2}$
\end{itemize}

\item \textbf{Feature engineering beats algorithm shopping}
\begin{itemize}
\item 1 hour of domain thinking > 3 hours of tuning
\end{itemize}

\item \textbf{Interpret your errors}
\begin{itemize}
\item 42 misclassified digits taught me more than 1,458 correct predictions
\end{itemize}

\item \textbf{Cross-platform reproducibility from day one}
\begin{itemize}
\item Use \texttt{os.path.join()} not hardcoded paths
\end{itemize}
\end{enumerate}

\end{frame}

% ============================================================================
% PART 3: THE BIG PICTURE (2 minutes)
% ============================================================================

\section{The Big Picture}

% Slide 9: When to Prefer Each Method
\begin{frame}{When to Prefer SVMs vs Tree-Based Methods}

\begin{columns}[T]
\begin{column}{0.48\textwidth}
\textbf{Prefer SVMs when:}
\begin{enumerate}
\item High-dimensional, sparse data\\
\small{(MNIST: 784 pixels)}
\item Smooth, geometric boundaries\\
\small{(Digit curves)}
\item Need confidence estimates\\
\small{(Decision function distance)}
\item Moderate training size\\
\small{(Thousands to tens of thousands)}
\end{enumerate}
\end{column}

\begin{column}{0.48\textwidth}
\textbf{Prefer Trees when:}
\begin{enumerate}
\item Interpretability paramount\\
\small{(Explain pixel > 0.5 decisions)}
\item Heterogeneous features\\
\small{(Mixed categorical + continuous)}
\item Missing values common\\
\small{(Surrogate splits)}
\item Massive training data\\
\small{(Millions of examples)}
\end{enumerate}
\end{column}
\end{columns}

\vspace{0.5cm}
\centering
\textbf{Match method to data structure, not personal preference}

\end{frame}

% Slide 10: Limitations
\begin{frame}{Limitations I Encountered}

\begin{columns}[T]
\begin{column}{0.48\textwidth}
\textbf{SVM Limitations:}
\begin{itemize}
\item Hyperparameter sensitivity\\
\small{Grid search $\gamma \in [10^{-4}, 10^{-1}]$}
\item Kernel selection = art\\
\small{RBF vs polynomial: trial/error}
\item Computational cost at scale\\
\small{Full MNIST prohibitive}
\item No built-in feature importance
\end{itemize}
\end{column}

\begin{column}{0.48\textwidth}
\textbf{Tree Limitations:}
\begin{itemize}
\item High-dim inefficiency\\
\small{Even with PCA reduction}
\item Extrapolation failure\\
\small{Constant predictions per leaf}
\item Instability\\
\small{Single tree: 92.2\%}\\
\small{Needed ensembles}
\end{itemize}
\end{column}
\end{columns}

\vspace{0.5cm}
\textbf{Both share:}
\begin{itemize}
\item Difficulty with online learning (must retrain from scratch)
\item Lack of theoretical guarantees on finite samples
\end{itemize}

\end{frame}

% Slide 11: Transformation
\begin{frame}{My Transformation: Algorithm-Centric → Data-Centric}

\begin{center}
\begin{tabular}{p{0.4\textwidth} p{0.4\textwidth}}
\toprule
\textbf{Before This Assignment} & \textbf{After This Assignment} \\
\midrule
"Complex data needs complex models" & "What is my data's structure?" \\
\midrule
Reach for ensembles by default & Fit simple baselines first \\
\midrule
Trust metrics alone & Validate with diagnostics \\
\midrule
Algorithm shopping & Feature engineering \\
\midrule
Count errors & Interpret errors \\
\bottomrule
\end{tabular}
\end{center}

\vspace{0.5cm}
\textbf{New Workflow:}
\begin{enumerate}
\item Understand data generating process
\item Fit simple baselines (linear + engineered features)
\item Diagnose failures (bias vs variance)
\item Match method to structure (high-dim + smooth → kernels)
\item Validate rigorously (diagnostics + algebraic verification)
\end{enumerate}

\end{frame}

% Slide 12: Key Takeaway
\begin{frame}{The Assignment's Greatest Gift}

\begin{center}
\Huge
\textcolor{blue}{Confidence to choose\\
simplicity when appropriate}
\end{center}

\vspace{1cm}

\Large
Model selection isn't about using the fanciest method.

\vspace{0.3cm}

It's about finding the \textbf{minimal complexity}\\
that captures the \textbf{true data structure}.

\vspace{0.5cm}

\normalsize
\begin{itemize}
\item Exercise 2: Linear data → Linear SVR optimal
\item Exercise 3: Curved boundaries → RBF SVM optimal
\end{itemize}

\end{frame}

% Closing slide
\begin{frame}
\begin{center}
\Huge Thank You

\vspace{1cm}

\Large Questions?

\vspace{1cm}

\normalsize
\textit{From Theory to Practice:\\
My Journey with Learning Machines}
\end{center}
\end{frame}

\end{document}
